% from ega4-21c-fr.tex
\oldpage[IV-4]{313}
peut, par un changement de base $\operatorname{Spec}(\mathcal O_{S,s})\to S$, se
ramener au cas où $S$ est un schéma local de point fermé $s$ et de point générique
$\eta$, et il faut montrer que, pour que $f_s$ soit un isomorphisme, il faut et il suffit
que $f_\eta$ le soit. Or, la suffisance de cette condition résulte de (21.12.15). Pour
montrer qu'elle est nécessaire, on raisonne comme dans la démonstration de (21.12.15)
(en remarquant, avec les mêmes notations, que la projection du point fermé de $S$ est le
point fermé de $S_0$) pour se ramener au cas où en outre $S$ est \emph{noethérien};
mais alors la conclusion résulte de $(\mathbf{III}, 4.6.7$, (ii)).
\end{proof}

\begin{remarks}[21.12.17]
\label{IV.21.12.17}
(i) La conclusion de la proposition (21.12.15) n'est plus valable si on élimine
l'hypothèse que les fibres $X_s$ sont irréductibles. Prenons en effet
$S=\operatorname{Spec}(A)$, où $A$ est un anneau de valuation discrète,
$Y=\mathbb P^1_S$, qui est propre et lisse sur $S$ (17.3.9). Notons encore $s$ et
$\eta$ le point fermé et le point générique de $S$; soit $z$ un point fermé de $Y_s$,
par exemple un point rationnel sur $k=k(s)$, et soit $X$ le $Y$-préschéma obtenu en
faisant éclater le point $z$. Comme l'anneau de polynômes $A[T]$ est régulier
$(\mathbf 0, 17.3.7)$ et de dimension $2$, on voit comme dans la démonstration de
(15.1.1.6) que, si $f:X\to Y$ est le morphisme structural, $f^{-1}(z)$ est isomorphe
à $\mathbb P^1_k$, et d'autre part, le complémentaire de $f^{-1}(z)$ dans $X_s$ est
isomorphe à $Y_s-\{z\}$, donc $Z_1=f^{-1}(z)$ et $Z_2$, adhérence dans $X_s$ du
complémentaire de $Z_1$, sont les deux composantes irréductibles de $X_s$. Par
ailleurs, $f$ est évidemment propre et $g=h\circ f$ est plat, car $X$ est intègre
$(\mathbf{II}, 8.1.4)$ et pour tout ouvert affine $U$ de $X$, l'homomorphisme
$A\to\Gamma(U,\mathcal O_X)$ est injectif $(\mathbf I, 1.2.7)$, donc
$\Gamma(U,\mathcal O_X)$ est un $A$-module sans torsion, et par suite plat puisque $A$
est un anneau de valuation discrète $(\mathbf 0_{\mathrm I}, 6.3.4)$. Il est clair que
$f_\eta:X_\eta\to Y_\eta$ est un isomorphisme, bien que
$f_s:X_s\to Y_s$ n'en soit pas un.

(ii) La conclusion de (21.12.15) tombe aussi en défaut si, dans les hypothèses, on
supprime l'hypothèse que $f$ est \emph{propre}. Il suffit pour le voir de définir $S$ et
$Y$ comme dans (i), mais de remplacer $X$ par le complémentaire $X'$ de $Z_2$ dans le
préschéma $X$ défini dans (i), $f$ par la restriction $f':X'\to Y$ de $f$; le morphisme
$g':X'\to S$, restriction de $g$, est encore plat, et cette fois $X'_s$ est
géométriquement irréductible; $X'_s$ est d'ailleurs isomorphe au complémentaire d'un
point fermé dans $\mathbb P^1_k$, donc est un schéma affine isomorphe à $\mathbb A^1_k$;
comme son image dans $Y_s$ est réduite au point fermé $z$, $f'$ n'est pas propre
$(\mathbf{III}, 4.4.2)$ et $f'_s$ n'est pas un isomorphisme, bien que $f'_\eta$ en soit
un.

(iii) Il se peut que l'énoncé de la proposition (21.12.15) reste valable quand on y
remplace le mot « isomorphisme » par « morphisme étale » (cf. (21.12.14, (v))). Il en
sera alors encore de même de (21.12.16).
\end{remarks}
\subsection{Couples parafactoriels. Anneaux locaux parafactoriels}
\EGAReaderTarget{ega4.c04.subsec.21-13}{subsection:IV.21.13}
\label{IV.21.13}

\begin{definition}[21.13.1]
\EGAReaderTarget{ega4.c04.def.21-13-1}{IV.21.13.1}
Soient $X$ un espace annelé, $Y$ une partie fermée de $X$, $U=X-Y$.
On dit que le couple $(X,Y)$ est \emph{parafactoriel} si, pour tout ouvert $V$ de $X$,
le foncteur restriction
$\mathscr L\mapsto\mathscr L|(U\cap V)$, de la catégorie des
$\mathcal O_V$-Modules inversibles dans celle des $\mathcal O_{U\cap V}$-Modules
inversibles, est une \emph{équivalence de catégories}.
\end{definition}

Dire que le couple $(X,Y)$ est parafactoriel signifie donc que, pour tout ouvert $V$
de $X$, les conditions suivantes sont vérifiées :
\begin{enumerate}[label=\arabic*\textsuperscript{o}]
  \item le foncteur $\mathscr L\mapsto\mathscr L|(U\cap V)$ est \emph{pleinement
  fidèle}, autrement dit, pour deux $\mathcal O_V$-Modules inversibles $\mathscr L$,
  $\mathscr L'$, l'application de restriction
  \[
    \operatorname{Hom}_{\mathcal O_V}(\mathscr L,\mathscr L')
    \longrightarrow
    \operatorname{Hom}_{\mathcal O_{U\cap V}}
    (\mathscr L|(U\cap V),\mathscr L'|(U\cap V))
  \]
  est \emph{bijective};
  \item le foncteur $\mathscr L\mapsto\mathscr L|(U\cap V)$ est
  \emph{essentiellement surjectif}, autrement dit, pour tout
  $\mathcal O_{U\cap V}$-Module inversible $\mathscr L_0$, il existe un
  $\mathcal O_V$-Module inversible $\mathscr L$ tel que $\mathscr L_0$ soit
  